% =====================================================================
%  ch00a.tex —— 引言（原书 p.512–517）
% =====================================================================
% [原文 p.512]
\chapter*{引言}
\addcontentsline{toc}{chapter}{引言}
\markboth{引言}{引言}

代数几何研究的主要对象是代数簇（或概形）以及它们之间的态射。每个代数簇 $X$ 都带有赋环拓扑空间的结构，因而它既有拓扑（通常是 Zariski 拓扑），又有正则函数环层 $\cO_X$。

在很大程度上，对代数簇的研究就是对其上诸层的研究。由于该空间是赋环的，自然的层就是其上的 $\cO_X$-模层，其中拟凝聚层与凝聚层以其代数本性而显得突出。回忆：一个 $\cO_X$-模层称为\emph{拟凝聚}的，如果它局部地可表示为一族自由层之间同态的余核；若这些自由层还是有限秩的，则称它为\emph{凝聚}的。（簇上的局部自由层与向量丛一一对应，因此下文不加区分地使用这两个术语。）

于是，每个代数簇 $X$ 都对应着两个阿贝尔范畴：凝聚层范畴 $\coh(X)$ 与拟凝聚层范畴 $\Qcoh(X)$。簇之间的态射在这两个阿贝尔范畴之间诱导出逆像函子与直接像函子。

然而这些函子并不正合，也就是说，它们不把正合列变为正合列。这使得在阿贝尔范畴以及其间非正合函子的框架下工作遇到相当大的困难。为保持函子性，Cartan 与 Eilenberg [11] 引入了导出函子的概念，对非正合函子作出必要的修正。这一技术由 Grothendieck 在 [15] 中加以发展，随后导致了新概念的引入：导出范畴，以及导出范畴之间的导出函子。

与阿贝尔范畴不同，导出范畴没有短正合列，态射的核与余核也没有定义。不过，导出范畴具有某种内部结构，Verdier 把它形式化为三角范畴的概念 [44]。

从阿贝尔范畴过渡到它们的导出范畴，使我们能够解决许多由研究自然函子而产生的困难问题。最早的一些例子中，可以提到整体相交理论的建立与黎曼--罗赫定理的证明。这些由 Grothendieck 及其合作者 [41] 得到的结果，正是由于引入了完美复形的三角范畴才成为可能。

另一个例子与 perverse sheaves 的引入有关，也与具有正则奇点的全纯模（holonomic modules）与可构造层之间的黎曼--希尔伯特对应的建立有关（见 [3]、[23]）；这一对应只有在运用三角范畴的概念与技术之后才成为可能。

许多与簇的研究有关的问题，都要求研究并描述其上的凝聚层导出范畴。在最简单的情形，即簇是一个点或一条光滑曲线时，凝聚层导出范畴中的每个对象都同构于某一族凝聚层的直和再作适当的平移；也就是说，每个 $A\in\Db(\coh X)$ 都同构于 $\bigoplus_{i=1}^{k}\cF_i[n_i]$，其中 $\cF_i$ 是凝聚层。这些例子反映了如下事实：在这些情形下，阿贝尔范畴的同调维数不超过 $1$。

% [原文 p.513]
然而，对于维数更高的簇，存在这样的复形：它在导出范畴中并不同构于其各阶上同调的直和。因此，描述维数大于 $1$ 的簇的导出范畴是一个困难而有趣的问题。这一方向上最初的几步由 [4] 与 [2] 迈出，它们描述了射影空间上凝聚层的导出范畴，随后使这一技术得以应用于 $\PP^2$ 与 $\PP^3$ 上向量丛的模空间的研究。特别地，这些论文表明，射影空间上凝聚层的导出范畴 $\Db(\coh\PP^n)$ 等价于有限维代数
\begin{equation*}
A=\End\Bigl(\bigoplus_{i=0}^{i=n}\cO(i)\Bigr)
\end{equation*}
上有限维模的导出范畴。此后这一方法日趋完善，二次超曲面与旗簇上凝聚层的导出范畴也得到了描述（[20]--[22]）。

引入例外族与半正交分解的概念，使人得以表述描述凝聚层导出范畴的新原则 [5]、[6]。结果表明，完备例外族的存在总实现凝聚层导出范畴与给定例外族的自同态代数（它是有限维代数）上有限维模的导出范畴之间的等价 [5]。半正交分解的概念则使我们能用被爆破簇以及沿之爆破的子簇的导出范畴，来描述爆破的导出范畴 [34]。

然而，对许多类型的簇，导出范畴是无法描述的。尽管如此，可以大致提出如下自然的问题：从簇过渡到其凝聚层导出范畴时，究竟保留了多少信息？事实上，在这一对应之下保留了“几乎全部”信息。在许多情形下，甚至可以从凝聚层导出范畴还原出簇本身，例如当典范层（或反典范层）丰富时 [8]。

不过，对某些类型的簇，仍然能找到两个不同的簇具有等价的凝聚层导出范畴的例子。两个不同簇具有等价凝聚层导出范畴的第一个例子由 Mukai 找到 [29]。他证明，对每个阿贝尔簇与其对偶簇，这都会发生。我们在 [38] 中推广了这一构造：对任一阿贝尔簇，我们引入了一整类阿贝尔簇，它们全都有相同的凝聚层导出范畴。一方面，这些例子表明存在具有等价凝聚层导出范畴的簇；另一方面，每一类具有等价凝聚层导出范畴的簇都是“小”的（在所有例子中它都是有限的）。

为得到具有等价凝聚层导出范畴的簇的完全分类，我们需要描述其间的函子与等价。结果表明，等价总是几何性的，也就是说，它们由两个簇的积上的某些层复形所表示。

% [原文 p.514]
我们来解释这句话的意思。以下用 $\Db(X)$ 记 $X$ 上凝聚层的有界导出范畴。光滑完备代数簇之间的任一态射 $f:X\to Y$ 在其凝聚层有界导出范畴之间诱导出两个正合函子：直接像函子 $\Rf f_*:\Db(X)\to\Db(Y)$，以及逆像函子 $\Lf f^*:\Db(Y)\to\Db(X)$，后者是前者的左伴随。此外，每个对象 $\cE\in\Db(X)$ 都定义一个正合张量积函子 $\otL\cE:\Db(X)\to\Db(X)$。利用这些标准的导出函子，可以引入 $\Db(X)$ 与 $\Db(Y)$ 之间一大类新的正合函子。

设 $X$ 与 $Y$ 是域 $k$ 上两个光滑完备簇。考虑积 $X\times Y$，并记
\begin{equation*}
X\xleftarrow{\ p\ }X\times Y\xrightarrow{\ q\ }Y
\end{equation*}
为 $X\times Y$ 到 $X$ 与到 $Y$ 的投影。每个对象 $\cE\in\Db(X\times Y)$ 都定义一个从导出范畴 $\Db(X)$ 到导出范畴 $\Db(Y)$ 的正合函子 $\Phi_{\cE}$，它由
\begin{equation}
\Phi_{\cE}(\,\cdot\,):=\Rb q_*\bigl(\cE\otL p^*(\,\cdot\,)\bigr) \tag{1}
\end{equation}
给出。这类函子都有左伴随与右伴随。

% [Tutor Note: 原书式 (1) 与 §2.1 以下各处对全导出函子的记号并不统一：本处印作 $\mathbf{R}^{\bullet}q_*$，后文（如 §2.2 定理 2.2.9 前的函子）印作 $\mathbf{R}\pi_*$。本译本一律记 $\mathbf{R}q_*$（必要时保留原书的 $\mathbf{R}^{\bullet}$），语义无别，仅记号从简。]

于是，每个光滑完备代数簇都可以配上一个凝聚层导出范畴；而两个这样的簇的积上的每个对象 $\cE\in\Db(X\times Y)$，都可以配上一个从三角范畴 $\Db(X)$ 到三角范畴 $\Db(Y)$ 的正合函子 $\Phi_{\cE}$。本文即致力于研究这一对应。

研究凝聚层导出范畴时，最先出现的问题之一是：这些范畴之间的每个函子都能由积上的一个对象表示吗？也就是说，它是否都具有 (1) 的形式？在第~3 章中，当该函子是等价时，我们对此问题给出肯定的回答。

另外两个核心问题如下：
\begin{enumerate}[label=\arabic*),leftmargin=2.4em,itemsep=2pt,topsep=2pt]
\item 两个不同光滑完备簇上的凝聚层导出范畴，何时作为三角范畴等价？
\item 给定簇 $X$ 上凝聚层导出范畴的正合自等价群是什么？
\end{enumerate}

这方面的一些结果已经知道。当簇具有丰富的典范层或反典范层时，上述问题已有完全的答案：在 [8] 中我们证明，具有丰富典范层（或反典范层）的光滑射影簇 $X$ 可由其凝聚层导出范畴 $\Db(X)$ 还原；此外，[8] 还给出了还原 $X$ 的具体构造。对这类簇，正合自等价群也可以描述。

下面我们说明本文的内容与结构。本文汇集的大部分结果都以某种形式见于 [7]、[8]、[34]、[35]、[37]。

第~1 章汇集预备性质的材料。我们首先给出三角范畴的定义，并回忆三角范畴之间正合函子的概念、三角范畴关于全子范畴的局部化，以及局部化三角范畴的导出函子的一般定义。此后我们定义阿贝尔范畴的同伦范畴与导出范畴，并讨论概形上凝聚层与拟凝聚层导出范畴的性质以及这些范畴之间的函子。

% [原文 p.515]
在第~2 章中，我们引入光滑完备代数簇上凝聚层有界导出范畴之间可由积上的对象表示的那一类函子，并描述其主要性质。利用第~3 章的结果，我们证明：若两个光滑射影簇 $X$ 与 $Y$ 的导出范畴等价，则如下公式所定义的分次代数 $\mathrm{HA}(X)$ 与 $\mathrm{HA}(Y)$ 之间存在同构：
\begin{equation*}
\mathrm{HA}(X)=\bigoplus_{i,k}\mathrm{HA}^{i,k}(X)
=\bigoplus_{i,k}\bigoplus_{p+q=i}H^{p}\Bigl(X,\bigwedge\nolimits^{q}T_X\otimes\omega_X^{k}\Bigr),
\end{equation*}
其中 $T_X$ 是 $X$ 的切丛，$\omega_X$ 是 $X$ 的典范丛（定理 2.1.8 与推论 2.1.10）。

在第~2 章第二节中，我们给出一整类具有等价凝聚层导出范畴的簇对。这些例子的有趣之处在于，所出现的簇彼此双有理同构（但一般不同构），并由一个称为 flop 的双有理变换相联系。特别地，这些例子表明，在由 $D(X)$ 还原 $X$ 的定理中，典范（或反典范）类丰富的条件不能减弱。

设 $Y$ 是光滑完备代数簇 $X$ 中光滑嵌入的闭子簇，满足 $Y\cong\PP^{k}$，且正规丛 $N_{X/Y}\cong\cO_Y(-1)^{\oplus(l+1)}$。我们假设 $l\leqslant k$，并用 $\tilde X$ 记 $X$ 沿中心 $Y$ 的爆破。此时例外除子 $\tilde Y$ 同构于射影空间的积 $\PP^{k}\times\PP^{l}$。存在 $\tilde X$ 的一个收缩，使 $\tilde Y$ 投影到第二个因子 $\PP^{l}$。考虑投影图
\begin{equation*}
X\xleftarrow{\ \pi\ }\tilde X\xrightarrow{\ \pi_+\ }X^{+}.
\end{equation*}
双有理映射 $\mathrm{fl}:X\dashrightarrow X^{+}$ 是 flip 或 flop 的最简单例子；当 $l<k$ 时它是 flip，当 $l=k$ 时它是 flop。

% [Tutor Note: 原书此处只写 “the simplest example of a flip or flop”，未定义 flip/flop 的一般概念；本译本保留英文原词，flip 首次出现时译作“翻转”并附原文，flop 全篇保留英文（见术语表）。]

本节的主定理把簇 $X$ 与 $X^{+}$ 上凝聚层的导出范畴联系起来。它断言：对 $\tilde X$ 上任一线丛 $L$，函子
\begin{equation*}
\Rf\pi_*\bigl(\Lf\pi_+^*(\,\cdot\,)\otimes L\bigr):\Db(X^{+})\to\Db(X)
\end{equation*}
是全忠实的；而当 $k=l$ 时，这个函子是等价。

第~3 章是全文的核心。它要证明：光滑射影簇上凝聚层导出范畴之间的每个等价，都由积上的一个对象表示。这一断言使我们能够描述凝聚层导出范畴之间的等价，并回答两个不同簇何时具有等价的凝聚层导出范畴这一问题。

事实上，本章证明一个更一般的断言：光滑射影簇上凝聚层有界导出范畴之间的任一函子，只要它是全忠实的并且具有伴随函子，就都可以由这些簇的积上的一个对象 $E$ 表示；也就是说，它同构于由规则 (1) 定义的函子 $\Phi_E$。此外，表示它的对象 $E$ 在同构意义下唯一确定（定理 3.2.1）。

% [原文 p.516]
在第~4 章中，我们研究 K3 曲面上凝聚层的导出范畴。对任一 K3 曲面 $S$，其上同调格 $H^{*}(S,\Z)$ 带有由
\begin{equation*}
(u,u')=r\cdot s'+s\cdot r'-\alpha\cdot\alpha'\in H^{4}(S,\Z)\cong\Z
\end{equation*}
给出的对称双线性型，其中 $u=(r,\alpha,s)$、$u'=(r',\alpha',s')$ 是 $H^{0}(S,\Z)\oplus H^{2}(S,\Z)\oplus H^{4}(S,\Z)$ 中的任意元素。带有该双线性型 $(\,\cdot\,,\,\cdot\,)$ 的上同调格 $H^{*}(S,\Z)$ 称为 Mukai 格，记作 $\tilde H(S,\Z)$。

格 $\tilde H(S,\Z)$ 带有自然的霍奇结构。在目前情形下，所谓霍奇结构，是指我们在空间 $\tilde H(S,\C)$ 中固定一个特殊的一维子空间 $H^{2,0}(S)$。若两个 K3 曲面 $S_1$ 与 $S_2$ 的 Mukai 格之间存在一个把一维子空间 $H^{2,0}(S_1)$ 映到 $H^{2,0}(S_2)$ 的等距，则称它们的 Mukai 格是霍奇等距的。

本章的主定理（定理 4.2.1）断言：域 $\C$ 上两个 K3 曲面 $S_1$、$S_2$ 上凝聚层的导出范畴 $\Db(S_1)$ 与 $\Db(S_2)$ 作为三角范畴等价，当且仅当它们的 Mukai 格之间存在霍奇等距 $f:\tilde H(S_1,\Z)\xrightarrow{\ \sim\ }\tilde H(S_2,\Z)$。此定理还有一个用超越闭链格表述的版本（定理 4.2.4）。

鉴于 K3 曲面的 Torelli 定理 [39]、[27]——它说 K3 曲面可由其第二上同调上的霍奇结构还原——我们得到了用霍奇结构表述的、关于两个 K3 曲面上凝聚层导出范畴何时等价的答案。

在第~5 章中，我们研究阿贝尔簇上凝聚层的导出范畴及其自等价群。设 $A$ 是阿贝尔簇，$\hat A$ 是对偶阿贝尔簇。如 [29] 所证，凝聚层的导出范畴 $\Db(A)$ 与 $\Db(\hat A)$ 等价，而这个称为 Fourier--Mukai 变换的等价，可由积 $A\times\hat A$ 上的 Poincar\'e 线丛 $P_A$ 按规则 (1) 给出：$\cF\mapsto \Rb p_{2*}\bigl(P_A\otimes p_1^*(\cF)\bigr)$。

Mukai 的这一构造在 [38] 中作了如下推广。考虑两个阿贝尔簇 $A$ 与 $B$，以及阿贝尔簇 $A\times\hat A$ 与 $B\times\hat B$ 之间的一个同构 $f$。把 $f$ 写成矩阵形式
\begin{equation*}
f=\begin{pmatrix} x & y\\ z & w\end{pmatrix},
\end{equation*}
其中 $x$ 表示从 $A$ 到 $B$ 的同态，$y$ 表示从 $\hat A$ 到 $B$ 的同态，其余类推。我们说同构 $f$ 是\emph{等距的}，如果它的逆具有形式
\begin{equation*}
f^{-1}=\begin{pmatrix} \hat w & -\hat y\\ -\hat z & \hat x\end{pmatrix}.
\end{equation*}
我们定义 $U(A\times\hat A,\,B\times\hat B)$ 为等距同构 $f$ 的集合。若 $B=A$，则将此集合记作 $U(A\times\hat A)$；注意它是 $\Aut(A\times\hat A)$ 的子群。

我们在 [38] 中证明：若代数闭域上两个阿贝尔簇 $A$、$B$ 的 $A\times\hat A$ 与 $B\times\hat B$ 之间存在等距同构，则凝聚层的导出范畴 $\Db(A)$ 与 $\Db(B)$ 等价。

% [原文 p.517]
在第~5 章中我们证明：在特征零的代数闭域上，这些条件是等价的；也就是说，$\Db(A)$ 与 $\Db(B)$ 等价，当且仅当存在从 $A\times\hat A$ 到 $B\times\hat B$ 的等距同构。事实上，“仅当”这一部分对任意域都成立（定理 5.1.16）。作为推论我们看到，对给定的阿贝尔簇 $A$，只有有限多个互不同构的阿贝尔簇，其导出范畴与 $\Db(A)$ 等价（推论 5.1.17）。

用积上的对象表示等价，我们构造了一个从 $\Db(A)$ 与 $\Db(B)$ 之间所有正合等价构成的集合，到 $A\times\hat A$ 与 $B\times\hat B$ 之间等距同构集合的映射。随后我们证明这个映射是函子性的（命题 5.1.12）。特别地，我们得到从 $\Db(A)$ 的正合自等价群到 $A\times\hat A$ 的等距自同构群 $U(A\times\hat A)$ 的一个同态。

在 \S\,5.2 中我们描述该同态的核，它同构于 $\Z$ 与 $A\times\hat A$ 的 $k$-值点群之直和（命题 5.2.3）。从技术上说，这一描述基于如下事实：定义等价的阿贝尔簇积上的那个对象，在导出范畴中相差一个平移的意义下，实际上是一个层（命题 5.2.2）。

在最后的 \S\,5.3 中，在基域代数闭且 $\operatorname{char}(k)=0$ 的假设下，我们对 [38] 中的断言再给出一个证明；这个证明用到 [30] 中关于阿贝尔簇上半齐次丛的结果。特别地，我们把自等价群描述为一个正合列
\begin{equation*}
0\longrightarrow \Z\oplus(A\times\hat A)_k\longrightarrow \Auteq \Db(A)\longrightarrow U(A\times\hat A)\longrightarrow 1.
\end{equation*}
